\section{Discussion}\label{sec:discussion}

\subsection{What the Results Mean for the Claim}\label{sec:disc-claim}

The question was whether one listener can serve every protocol at no measurable cost. For this
server, under load of one protocol at a time, the answer is yes in \CostClaimCount{} of the
configurations tested. In each, churn throughput, keep-alive throughput and the median time to
first byte stayed within the margin of \MarginLow{} to \MarginHigh{} of dedicated mode. These
configurations span two hosts, three backends and three protocols, HTTP/1.1, TLS and MQTT, and
nothing outside them.

Equivalence within the margin is the claim; a cost of zero is not. All C1 and C3 medians on L lie
on dedicated mode's side of one, and so do the server's CPU ratios of those cells. No test of a
nonzero cost was planned, and the interval of C1's TLS cell on io\_uring contains one.

The cost claim does not hold for h2c on L, which fails C3 on both backends. On epoll the analysis
reports a measured cost, which appears as latency at a fixed load, since h2c's churn and keep-alive
cells pass. On W the claim is open for h2c and MQTT: their churn measured the generator, and their
open-loop cells were not resolved. Two further results bound what one listener can promise. The
mixed-protocol cell on io\_uring had no valid session, in both modes. That is a limitation of the
server's io\_uring backend, not a cost of detection, but the cost claims cannot be extended to mixed
load on io\_uring. And sslh-ev held a smaller counted footprint per pending connection than the
server, in both cases and on both backends. So did Netty, a descriptive cell, in the
partial-ClientHello case.

B1 and B2 hold where they were tested. The detector decided every hard case as the frozen table
specifies, on every backend and in every mode, and no timed event broke its bounds. B3 holds against
nginx, HAProxy and Envoy in both cases and hyper-util in the silent case. By B3's test each held at
least \FootprintMargin{} times the server's counted footprint per pending connection, as
configured. The differences range from \BThreeSilentHyperUtilIouringDMedian{} to
\BThreePartialHelloEnvoyEpollDMedian{} bytes, with the exclusions and residuals of WL7
(Section~\ref{sec:res-robust}). B3 is lost against sslh-ev. The mechanisms mostly behave as
designed. Replay beat peek in \CountPassMOne{} of \CountHolmMOne{} cells, and in-process dispatch
used much less CPU than relay for HTTP/1.1. The server's relay beat every proxy it could be tested
against, which excludes sslh-ev. These are results for this server's design, not for one-port
listening in general.

\subsection{Hypotheses for Later Work}\label{sec:disc-hyp}

The runs measured where the claim fails but not why. The following are hypotheses only; no run
of this paper tested them.

\runin{h2c at a fixed load on L.} The extra latency does not track extra work. The CPU ratio of C3's
h2c cell on epoll, \SecWLFourCThreeLEpollHTwoCCpuMedian, is smaller than HTTP/1.1's,
\SecWLFourCThreeLEpollHttpCpuMedian, whose C3 cell passes. On io\_uring the two are about equal,
\SecWLFourCThreeLIouringHTwoCCpuMedian{} and \SecWLFourCThreeLIouringHttpCpuMedian. So we suspect a
wait rather than work: a step on h2c's path in one-port mode that delays the server's first byte
without using the CPU. A trace of each connection's events in both modes, at the C3 rate, would
locate it.

\runin{The mixed cell on io\_uring.} The backend's ring holds \RingBuffers{} free buffers per worker,
and the cell's background keeps as many keep-alive connections busy. We suspect that the background's
receives, completing together, empty the ring before the worker gives it new buffers. The churn's
receives would then fail and be posted again, and its probe would time out. The server counts such
retries, and this paper reports no value of that counter. The epoll backend, which reads into each
connection's own buffer, served the same cell. A larger ring is a compiled change after the code
freeze, so this paper does not make it. It is planned as an exploratory follow-up, labelled as such,
and the sizing of the ring is later work.

\runin{sslh-ev's memory.} Both systems keep a pending connection near the same cost, so neither can
reach the margin against the other. Two parts may explain sslh-ev's lead. In the
partial-ClientHello case its $U$ per pending connection, \BThreePartialHelloSslhEvEpollCompetitorU{}
bytes on epoll, is below the server's, \BThreePartialHelloSslhEvEpollServerU. We suspect the buffer
each holds for the bytes received so far. On io\_uring the server's $K_s$ exceeded \Kbase{} in most
cells, while every competitor's fell below it. In the silent case $K_s - \Kbase$ is
$\BThreeSilentSslhEvIouringServerKsNet$ bytes for the server against
$\BThreeSilentSslhEvIouringCompetitorKsNet$ for sslh-ev. We suspect kernel memory that the io\_uring
backend holds per connection. A server that held a smaller buffer until the first decision, and a
count of the slab objects per backend, could test both.

\runin{Two listeners on io\_uring.} A \texttt{SO\_REUSEPORT} group served
\SecTwoCoresSecTwoCoresCOneLIouringReuseportValueMedian{} times the shared listener's connections
per second on io\_uring, but about the same on epoll. We suspect that the two workers' multishot
accepts on one shared socket do not spread connections evenly between the workers. Counting each
worker's accepts would test it. That question belongs to the accept models, outside this paper.

\subsection{Threats to Validity}\label{sec:threats}

\runin{Scope.} Every claim is about one server, as built, on two hosts, over loopback. A network
between client and server adds latency and interrupt work that loopback lacks, and we expect both to
dilute the cost of detection further. The cost claims hold for load of one protocol at a time. The
mixed-protocol cell is secondary, and on io\_uring it had no valid session. Each claim is made per
protocol, backend and host, never pooled.

\runin{What the cells measure.} The TLS cells of the cost family have little power to show the cost
of detection, since the handshake dominates each connection. In the SSH cells the two arms exchange
messages in a different order. M1 on IOCP compares replay with a peek that would switch to replay at
its first undecided read, because Windows offers no low-water mark; in the runs no peek switched. On
W the open-loop cells of h2c and MQTT ran at rates taken from generator-bound churn, so they ran at
a smaller share of the server's capacity than open HTTP/1.1.

\runin{Engineering before the freeze.} The server was engineered until it won where it could, on
development data, before the code freeze. Those runs could not time one-port mode against dedicated
mode, and the revision log lists no pair of runs that would give such a ratio indirectly. The A/A
pilot that sized the cost family held no one-port data. So neither the margin nor \RC{} could be
tuned to a cost contrast. The relay and the footprint test were informed by development runs
against the competitors, and a reader should weigh that history.

\runin{The footprint test.} $W$ is the footprint of each system as configured, not of its detection
design alone. Two residuals of the socket-buffer accounting remain, and both sit with a system that
leaves bytes queued: nginx and Envoy in the partial-ClientHello case, and the server only in peek.
The co-tenants of the merged cache on L grow only when the kernel creates workqueue pools or chained
scatter-gather lists. The merged cache's growth is reported per arm in
Tables~\ref{tab:b3parts} and~\ref{tab:b3partsserver}, and in the silent case no system leaves a byte
queued. Against a system whose footprint is near a bare socket's, little room is left for the
server to reach the margin. sslh-ev is such a system, as development runs before the freeze had
suggested.
The
JVM and Go systems hold memory that their collectors manage, and their cells are descriptive.

\runin{Windows.} W's windows were validated without a frequency rule, since no counter passed its
test, so a change of clock speed within a W session would go undetected. W cannot count listen
overflows, so a refused connection shows only as a failed connect, which invalidates its window.
W's job-start check was lowered, so a job could start on a busier W than first planned; every
window rule still applied.

\runin{Cells that did not run.} Three outcomes were foreseen before the runs. M2's TLS cells were not
run, which narrows M2 to HTTP/1.1, and sslh-ev's M3 cells are untested, which narrows M3. TLS with
session resumption, a secondary cell, was not run and narrows nothing. Decided before the
confirmatory runs, W's churn of h2c and MQTT ran in the pilot only, and the two-core cell on IOCP
was not run. Not foreseen were the SSH cells of C2, which could not run, and the mixed cell on
io\_uring. Four descriptive outputs were produced after the analysis, and the system-call counts
and the relay's operations came from an untimed job run then (Section~\ref{sec:dev-not-produced}).

\runin{Competitors and instruments.} Each competitor runs the configuration its project documents
for the compared feature, and a better one may exist. HAProxy's client, server and connect timeouts
stay at its default, which is none. The runs use kernel \KernelL, while the pre-specification cites
the source at tag \KernelTag; each cited line is unchanged there, but the distribution's own patches
were not compared. A closed-loop cell measures throughput at \ConnSlots{} connections per server
core only.

\runin{Analysis.} Two hosts with different Python versions computed every decision equal
(Section~\ref{sec:dev-escaping}). B1's hard cases and B2 were read from the runners' rows rather than
computed by the analysis code, as a logged reading states (Section~\ref{sec:dev-runners}).
